%versi 3 (22-07-2020)
\chapter{Perancangan}
\label{chap:perancangan}

Bab ini membahas rancangan perangkat lunak berdasarkan hasil analisis pada Bab~\ref{chap:analisis}. Pada bab sebelumnya, fokus pembahasannya adalah masalah, kebutuhan sistem, dan gambaran sistem usulan. Pada bab ini, pembahasan mulai diarahkan ke cara sistem tersebut akan dibangun. Rancangan dibuat agar proses dari URL sampai menjadi visualisasi graf memiliki alur yang jelas dan setiap bagian program mempunyai tugas yang tidak saling tumpang tindih.

\section{Perancangan Arsitektur Sistem}
\label{sec:perancangan-arsitektur-sistem}

Arsitektur sistem dirancang dengan membagi proses utama ke dalam beberapa komponen. Pembagian ini dilakukan agar program lebih mudah dipahami, dikembangkan, dan diperbaiki. Dengan demikian, sistem tidak perlu mengerjakan semuanya sekaligus. Setiap komponen cukup menangani tugas tertentu, kemudian hasilnya diteruskan ke komponen berikutnya.

Proses dimulai dari komponen masukan URL yang menerima alamat halaman web dari pengguna. URL tersebut terlebih dahulu diperiksa agar sistem tidak langsung memproses masukan yang kosong atau memiliki format yang tidak sesuai. Setelah URL dinyatakan valid, komponen \textit{HTML Fetcher} menggunakan Jsoup untuk mengambil dokumen HTML dan membentuk objek \texttt{Document}. Objek ini menjadi sumber data utama untuk proses penelusuran DOM.

Komponen \textit{DOM Traversal} kemudian menelusuri elemen-elemen pada objek \texttt{Document} menggunakan metode DFS. Setiap elemen yang ditemukan dikirim ke \textit{Graph Builder} untuk dibentuk menjadi \textit{node}. Pada saat yang sama, hubungan langsung antara elemen induk dan elemen anak dicatat sebagai \textit{edge}. Agar elemen dengan nama \textit{tag} yang sama tidak tertukar, setiap \textit{node} diberikan ID unik, sedangkan nama \textit{tag} tetap digunakan sebagai label yang nantinya terlihat pada gambar graf.

Hasil pembentukan \textit{node} dan \textit{edge} diteruskan kepada \textit{DOT Generator}. Komponen ini menyusun data graf ke dalam format DOT yang dapat dibaca oleh Graphviz. Setelah itu, \textit{Graph Visualizer} meminta Graphviz mengubah data DOT menjadi gambar. Gambar tersebut akhirnya ditampilkan melalui bagian keluaran bersama informasi pendukung, seperti daftar \textit{node}, daftar \textit{edge}, dan teks DOT. Singkatnya, arsitektur sistem membentuk alur berurutan dari masukan URL, pengambilan HTML, pembentukan DOM, penelusuran elemen, penyusunan graf, sampai penampilan hasil visualisasi.

\section{Perancangan Alur Pengolahan DOM}
\label{sec:perancangan-alur-pengolahan-dom}

Alur pengolahan DOM dirancang agar seluruh elemen HTML dapat diperiksa tanpa kehilangan hubungan hierarkisnya. Sistem memulai proses dari elemen akar dokumen, lalu bergerak ke elemen anak paling awal. Jika elemen tersebut masih memiliki anak, penelusuran dilanjutkan lagi ke tingkat yang lebih dalam. Ketika sebuah elemen sudah tidak memiliki anak, sistem kembali ke elemen sebelumnya dan melanjutkan ke cabang berikutnya. Pola ini mengikuti cara kerja DFS yang telah dipilih pada tahap analisis.

Saat sebuah elemen pertama kali dikunjungi, sistem menambahkan satu nomor pada pencacah ID dan membentuk ID seperti \texttt{node1}, \texttt{node2}, dan seterusnya. Nama \textit{tag} elemen, misalnya \texttt{html}, \texttt{body}, atau \texttt{div}, disimpan sebagai label. Jika elemen tersebut memiliki induk, sistem juga menambahkan \textit{edge} dari ID induk menuju ID elemen yang sedang diproses. Dengan cara ini, susunan graf tetap mengikuti hubungan \textit{parent-child} pada DOM aslinya.

Setelah penelusuran selesai, seluruh deklarasi \textit{node} dan \textit{edge} digabungkan ke dalam blok \texttt{digraph}. Label yang mengandung karakter khusus perlu melalui proses \textit{escaping} agar tidak merusak sintaks DOT. Hasil DOT kemudian diberikan kepada Graphviz untuk dibuat menjadi gambar graf berarah. Kalau prosesnya berhasil, sistem menampilkan gambar dan detail data graf kepada pengguna. Kalau terjadi masalah, misalnya URL tidak dapat diakses, HTML gagal diproses, atau Graphviz tidak berhasil membuat gambar, sistem menampilkan pesan kesalahan sesuai tahap yang bermasalah.

Intinya, rancangan alur ini menjaga agar setiap elemen hanya dibuat sebagai satu \textit{node}, setiap hubungan langsung menghasilkan satu \textit{edge}, dan seluruh cabang DOM tetap ditelusuri. Dengan aturan tersebut, visualisasi yang dihasilkan diharapkan tetap sesuai dengan struktur DOM halaman web yang diterima oleh Jsoup.